\begin{summarybox}
\textbf{\textcolor{CSummary}{一句话记住它为什么成立}}

\textbf{内心到三边同样远}，所以连接内心与三个顶点后，
三块小三角形以三边为底时\textbf{高都是 $r$}；
三块面积相加就是原三角形面积。

\medskip
\textbf{\textcolor{CSummary}{必会的关系}}
\[
\boxed{ID=IE=IF=r},\qquad
\boxed{S=rp=\frac12r(a+b+c)},
\qquad \boxed{r=\frac{2S}{a+b+c}=\frac Sp}.
\]
\[
\boxed{S_{\triangle IBC}:S_{\triangle ICA}:S_{\triangle IAB}=a:b:c}.
\]
其中 $p=(a+b+c)/2$。
\[
\boxed{AE=AF=x=p-a},\quad
\boxed{BF=BD=y=p-b},\quad
\boxed{CD=CE=z=p-c}.
\]
切线段长沿三角形的边量取，等于半周长减去出发顶点的对边长。

\medskip
\textbf{\textcolor{CSummary}{考场上怎么用}}
\begin{enumerate}[leftmargin=*,itemsep=0.35em]
 \item \textbf{先识别：}题中有内心、内切圆、面积、周长或内心分出的面积比。
 \item \textbf{再选择：}已知 $S,L$ 求 $r$，用 $r=2S/L$；
 已知 $r,L$ 求 $S$，用 $S=rL/2$；三边已知求 $r$ 时先求面积。
 \item \textbf{能省则省：}只求三块面积比，直接用底边比，不求 $r$；
 直角三角形可用 $r=(u+v-w)/2$。
 \item \textbf{求切点分段：}三边已知，先算 $p$，再用 $x=p-a,y=p-b,z=p-c$，无需求 $r$。
\end{enumerate}

\medskip
\textbf{\textcolor{CSummary}{特殊三角形速查（先按理解说明建系）}}

等边三角形边长 $\ell$：$r=\sqrt3\ell/6$，$I=(\ell/2,r)$，$x=y=z=\ell/2$。

直角顶点在原点、两直角边沿正坐标轴：$I=(r,r)$，$r=(u+v-w)/2$。

等腰直角三角形的直角边为 $t$：$r=(2-\sqrt2)t/2$，$I=(r,r)$；
若给的是斜边 $w$，则 $r=(\sqrt2-1)w/2$。

\medskip
\textbf{\textcolor{CSummary}{最后检查}}

仅适用于真实的三角形；注意距离要\textbf{垂直测量}；
内心到\textbf{三边}等距而非到三个顶点等距；
看到 $p$ 记得它是\textbf{半周长}，不要多除或少乘一个 $2$。
\end{summarybox}
